\documentclass[aspectratio=169,11pt]{beamer}

\usetheme{Madrid}
\usecolortheme{default}
\usefonttheme{professionalfonts}
\setbeamertemplate{navigation symbols}{}
\setbeamertemplate{footline}[frame number]

\usepackage{amsmath, amssymb, booktabs}

\title{Mental Health, Stress, Workplace Productivity, And Worker Welfare}
\subtitle{Absenteeism, Presenteeism, Stigma, Job Quality, And Causal Limits}
\author{Special Topic 6: Labor Market, Health, and Population -- Week 4}
\date{}

\begin{document}

\begin{frame}
  \titlepage
\end{frame}

\begin{frame}{Central question}
\begin{itemize}
    \item How do mental health and stress affect productivity, job quality, retention, and welfare inside labor markets?
    \item Mental health is not an individual residual; it is a labor-market force affecting attendance, effort, search, supervision, safety, and treatment.
    \item Week 4 discipline: distinguish mental health affecting work, work affecting mental health, treatment and stigma, and measurement and welfare.
\end{itemize}
\end{frame}

\begin{frame}{Mental health as productivity and welfare object}
\small
\begin{tabular}{p{0.28\linewidth} p{0.32\linewidth} p{0.28\linewidth}}
\toprule
Object & Labor margin & Welfare interpretation \\
\midrule
Depression / anxiety & concentration, effort, social functioning & work may be more costly \\
Stress and pain & absences, errors, safety, fatigue & wage may miss burden \\
Substance use & attachment, reliability, search & dynamic scarring risk \\
Treatment & attendance, retention, recovery & access and timing matter \\
Job quality & autonomy, conflict, meaning & amenity and health object \\
\bottomrule
\end{tabular}
\vspace{0.25cm}
\begin{itemize}
    \item Wages and employment can miss absenteeism, presenteeism, stigma, safety risk, and diminished job quality.
\end{itemize}
\end{frame}

\begin{frame}{Productivity margins inside work}
\begin{itemize}
    \item Absenteeism: missed days, sick leave, disability leave, and episodic nonattendance.
    \item Presenteeism: reduced concentration, speed, reliability, interpersonal functioning, or task completion while at work.
    \item Performance and safety: output, errors, customer interaction, supervisor ratings, and injury risk.
    \item Retention: quits, job-to-job moves, delayed promotion, occupational downgrading, and exit.
    \item Firm policy: scheduling, supervision, leave, accommodations, employee assistance, and monitoring.
\end{itemize}
\end{frame}

\begin{frame}{Stigma, treatment, and disclosure}
\small
\begin{tabular}{p{0.27\linewidth} p{0.34\linewidth} p{0.27\linewidth}}
\toprule
Mechanism & Worker response & Labor outcome \\
\midrule
Stigma & hide symptoms or treatment & delayed support \\
Disclosure risk & avoid asking for help & no accommodation \\
Treatment friction & postpone care & absences or presenteeism \\
Manager beliefs & interpret distress as risk & promotion and retention effects \\
Career concerns & sort away from some jobs & constrained choice set \\
\bottomrule
\end{tabular}
\vspace{0.25cm}
\begin{itemize}
    \item Stigma is a labor-market mechanism because it changes information, help-seeking, accommodations, attendance, and careers.
\end{itemize}
\end{frame}

\begin{frame}{Labor conditions shape mental health}
\begin{itemize}
    \item Workload and time pressure can raise stress and reduce recovery time.
    \item Autonomy and control can buffer stress by changing pace, methods, scheduling, and task order.
    \item Supervision and monitoring can provide support or generate anxiety and quits.
    \item Conflict, harassment, customer abuse, and discrimination can make work a direct source of mental-health harm.
    \item Isolation, remote work, mission alignment, and meaning affect support, identity, and welfare.
\end{itemize}
\end{frame}

\begin{frame}{Measurement and causal challenges}
\small
\begin{tabular}{p{0.24\linewidth} p{0.32\linewidth} p{0.30\linewidth}}
\toprule
Measure & What it captures & Main caution \\
\midrule
PHQ-9, GHQ, CES-D, K6 & symptoms and distress & differential reporting \\
Prescriptions / claims & treatment contact & access and stigma \\
Sick leave / disability & work disruption & program and workplace rules \\
Absenteeism & missed work & leave access and norms \\
Presenteeism & on-the-job loss & hard to observe cleanly \\
Performance / quits & workplace outcomes & job quality and selection \\
\bottomrule
\end{tabular}
\vspace{0.25cm}
\begin{itemize}
    \item Identification threats: reverse causality, dynamic selection, diagnosis intensity, treatment selection, co-morbidity, SES confounding, and reporting.
\end{itemize}
\end{frame}

\begin{frame}{Frontier methods and research opportunities}
\begin{itemize}
    \item Treatment-access shocks: provider supply, coverage, copayments, employee assistance, and appointment access.
    \item Waiting-time shocks: delayed treatment as variation in recovery and labor attachment.
    \item Job-loss and workplace shocks: displacement, reorganization, monitoring, remote work, and pandemic-era changes.
    \item Policy changes: paid leave, flexibility, accommodation rules, disability rules, and mental-health coverage.
    \item Structural and job-search approaches: jointly model mental health, treatment, job choice, productivity, and welfare.
\end{itemize}
\end{frame}

\begin{frame}{Research Lab design}
\begin{itemize}
    \item Primary anchor: Bubonya, Cobb-Clark, and Wooden on mental health and productivity at work.
    \item Challenge anchor: job meaning, identity, and employee mental health.
    \item Reproduce: estimate a bounded synthetic profile linking distress to absenteeism, presenteeism, performance, quits, and welfare.
    \item Diagnose: separate productivity loss, job-quality effects, reporting, stigma, treatment selection, and reverse causality.
    \item Transfer: redesign the logic for treatment access, waiting time, job loss, workplace shocks, disclosure interventions, or structural search.
\end{itemize}
\end{frame}

\begin{frame}{Research design checklist}
\begin{enumerate}
    \item Name the latent object: distress, diagnosis, treatment, functioning, stigma, or workplace loss.
    \item Name the labor margin: attendance, presenteeism, performance, quits, promotion, safety, or welfare.
    \item State the direction: mental health to work, work to mental health, or treatment and stigma to both.
    \item Identify the main threat: reverse causality, selection, reporting, diagnosis intensity, co-morbidity, or SES.
    \item Say what wages and employment miss.
\end{enumerate}
\end{frame}

\begin{frame}{Bridge to Week 5}
\begin{itemize}
    \item Week 4 centers mental health, stress, stigma, job quality, productivity, retention, and worker welfare inside workplaces.
    \item Week 5 scales up to disease exposure, environmental health, demographic change, and labor-market adjustment.
    \item The common thread is health as a force shaping workers, firms, allocation, productivity, and welfare.
\end{itemize}
\end{frame}

\begin{frame}{Takeaways}
\begin{itemize}
    \item Mental health affects labor markets through productivity, attendance, search, retention, safety, treatment, disclosure, and welfare.
    \item Work conditions also produce mental-health outcomes through workload, autonomy, supervision, conflict, isolation, and meaning.
    \item Measurement is the core empirical problem: symptoms, treatment, claims, absences, and performance are different objects.
    \item The frontier is open because many estimates are correlational or structurally identified under strong assumptions.
\end{itemize}
\end{frame}

\end{document}
